% TOP_PROBLEM: {"problemId": "problem.top500-omission-w045049-exponential-algebraic-closedness", "canonicalTitle": "Zilber's Exponential-Algebraic Closedness Conjecture", "releaseRank": 441, "primaryDomain": "logic_foundations_set_theory", "primaryDomainLabel": "Logic, foundations & set theory", "displayStatus": "open_with_solved_subcases", "releaseStatus": "open_with_solved_subcases"}
% !TeX program = lualatex
% ATTEMPT_STATUS: unresolved
% Research continuation by GPT_6_Astra_Ultra, 27 September 2026.
\documentclass[11pt]{article}
\usepackage[margin=25mm]{geometry}
\usepackage{fontspec}
\setmainfont{Latin Modern Roman}
\usepackage{amsmath,amssymb,mathtools}
\usepackage{unicode-math}
\setmathfont{Latin Modern Math}
\usepackage{xurl,hyperref}
\hypersetup{colorlinks=true,urlcolor=blue}
\setcounter{secnumdepth}{0}
\setlength{\parindent}{0pt}
\setlength{\parskip}{5pt}
\setlength{\emergencystretch}{3em}
\begin{document}
\section*{\texorpdfstring{Zilber's Exponential-Algebraic Closedness Conjecture}{Zilber's Exponential-Algebraic Closedness Conjecture}}\label{441}
\subsection{Definitions and mathematical statement}
Work in ${\mathbb{C}}^n\times({\mathbb{C}}^*)^n$. A variety $V$ is free if its additive projection is not contained in a translate of a proper rational linear subspace and its multiplicative projection is not contained in a translate of a proper algebraic subtorus. For a linear subspace $L\subseteq{\mathbb{C}}^n$ defined over ${\mathbb{Q}}$, the coordinatewise exponential image $\exp(L)$ is an algebraic subtorus. Write $\pi_L$ for the quotient map to
\[
 ({\mathbb{C}}^n/L)\times(({\mathbb{C}}^*)^n/\exp(L)).
\]
Call $V$ rotund if $\dim\pi_L(V)\ge n-\dim L$ for every such $L$. The conjecture asks that every irreducible free rotund $V$ contain an exponential point:
\[
 \exists z\in{\mathbb{C}}^n:\quad (z_1,\ldots,z_n,e^{z_1},\ldots,e^{z_n})\in V.
\]
Only existence is requested, not genericity over a coefficient field or a prescribed transcendence degree.
\par\smallskip\textit{Formulation and concept references:} \href{https://link.springer.com/article/10.1007/s00029-023-00853-y}{[S1]}.

\subsection{Short English statement}
Must every algebraic variety with the specified freedom and dimension conditions meet the graph of complex exponentiation?
\subsection{Research attempt}
\textit{GPT-6 Astra Ultra; 27 September 2026. Status: UNRESOLVED.}

I test two genuinely different features of the hypothesis. First I prove
existence for a family in which the algebraic and exponential coordinates
are coupled, but separate coordinate by coordinate. Then I construct a
free variety with no exponential point when rotundity is dropped. This
both produces an affirmative subcase and shows why freedom alone cannot
replace the dimension inequalities. No transcendence conjecture or
genericity conclusion is used.

\paragraph{A coupled family with an elementary existence proof.}
Choose $a_i\in\mathbb C^*$ and $b_i\in\mathbb C$, and let
\[
 V=\{(z,y):y_i=a_i z_i+b_i\ (1\le i\le n)\}
 \subset\mathbb C^n\times(\mathbb C^*)^n.
                                                        \tag{1}
\]
This is an irreducible variety of dimension $n$, isomorphic to the
product of the affine lines with the points $-b_i/a_i$ removed. Its
additive projection is dense in $\mathbb C^n$, and its multiplicative
projection is the entire torus. Thus neither projection lies in one of
the forbidden proper translates. For every rational linear $L$, the
additive component of $\pi_L(V)$ is dense in $\mathbb C^n/L$.
Consequently $\dim\pi_L(V)\ge n-\dim L$, proving rotundity with all
the required quantifiers, not just for coordinate subspaces.

To find an exponential point, it suffices to solve $e^z=az+b$ for every
$a\ne0$. Set $w=-z-b/a$. The equation becomes
\[
 we^w=c,\qquad c=-e^{-b/a}/a\ne0.                     \tag{2}
\]
Here is a direct existence argument without assuming a special inverse
function. Fix a logarithm $\lambda$ of $c$, put
$L_k=\lambda+2\pi i k$, and take positive integers $k$ sufficiently
large. On the closed disk
\[
 D_k=\{w:|w-L_k|\le2\log|L_k|\}
\]
the principal logarithm is analytic: the disk lies strictly in the
upper half-plane. It also has $|w|\ge|L_k|/2$. For sufficiently large
$k$,
\[
 |\operatorname{Log}w|\le\log(2|L_k|)+\pi
 \le2\log|L_k|.
\]
Thus $T_k(w)=L_k-\operatorname{Log}w$ maps $D_k$ into itself.
The derivative bound $|T_k'(w)|\le2/|L_k|<1$, integrated along the
line segment in this convex disk, proves it is a contraction. The
fixed-point theorem gives $w_k+\operatorname{Log}w_k=L_k$, and
exponentiating gives (2). Distinct $k$ give distinct fixed points,
since the same principal logarithm is used in their equations.
Therefore each coordinate equation in (1) has infinitely many solutions,
and selecting one for each coordinate proves the required existence.

This calculation also checks degeneracy. The assumption $a_i\ne0$
ensures additive dominance. If $a_i=0$, the equation may still have
solutions when $b_i\ne0$, but the multiplicative projection has a
constant coordinate and the displayed family no longer satisfies the
same freedom hypothesis. I do not extend the family by discarding that
condition.

\paragraph{A free variety without an exponential point.}
The following construction deliberately fails rotundity and therefore
does not contradict the conjecture. Put $\alpha=\sqrt2$. The entire
function
\[
 h(s)=e^{\alpha s}-e^s-1
\]
is not identically zero, since $h(0)=-1$. Its zero set $Z$ is discrete
and countable, and excludes zero. For a fixed $s\in Z$, the set of
$\beta$ satisfying
\[
 e^{\beta s}=e^s+2                                      \tag{3}
\]
is countable if the right side is nonzero, and empty otherwise. Exclude
all these sets, as well as $\mathbb Q+\mathbb Q\sqrt2$, and choose
$\beta$ outside their countable union. Such a complex number exists.

Consider
\[
 W=\{(s,\alpha s,\beta s;t,t+1,t+2):
 s\in\mathbb C,\ t\notin\{0,-1,-2\}\}.
                                                        \tag{4}
\]
It is closed in the stated additive-times-torus ambient space, being
defined there by linear equations, and is irreducible of dimension two.
An additive rational relation constant on its projection must vanish at
$s=0$ and has coefficient $q_1+q_2\alpha+q_3\beta$. Our choice makes
that coefficient nonzero unless all $q_i$ vanish. Hence its additive
projection is free.

A multiplicative relation constant on its projection would say
\[
 t^{m_1}(t+1)^{m_2}(t+2)^{m_3}=c_0
 \quad(m_i\in\mathbb Z).
\]
As an identity of rational functions, its valuations at $0,-1,-2$
force respectively $m_1=m_2=m_3=0$. Every proper subtorus translate
imposes a nontrivial such character relation, so the multiplicative
projection is also free. If (4) contained an exponential point, its
first coordinate would give $t=e^s$, the second would give $s\in Z$,
and the third would give (3), which was excluded. It contains none.
Finally, taking $L=0$ in the rotundity condition would require
$\dim W\ge3$, whereas $\dim W=2$. The failed hypothesis is completely
visible.

\paragraph{Why these arguments do not extend automatically.}
In (1), the algebraic equations decouple, and a large logarithmic branch
parameter can be chosen independently in each coordinate. General
irreducible varieties may couple every coordinate, and the contraction
map used above then has neither a defined domain nor a controlled
Jacobian without further work. Rotundity controls algebraic dimensions
of quotient images; it is not a lower bound for the derivative of an
analytic parametrization. Turning it into such a bound would be an
additional theorem, not a consequence of the displayed inequalities.

The example (4) shows that a dimension-free argument based on the two
projections being free cannot be repaired merely by choosing more
logarithm branches. Those branches were explicitly exhausted by the
countable exclusions defining $\beta$. For a rotund variety, the needed
argument must use the quotient-dimension conditions to prevent this
kind of simultaneous avoidance.

\paragraph{Literature, verification, and outcome.}
Gallinaro's theorem [S1, E1] proves the conjecture for varieties splitting
as a linear additive space times an algebraic torus subvariety. The
general variety in the catalogue need not split in that way. The proof
for (1) above is elementary and does not assert an extension of that
theorem to arbitrary coupled equations. Verification consisted of
checking every quotient condition for (1), the logarithm domain and
contraction estimates, and all freedom and failure conditions for (4).
No finite numerical root search is used as evidence of nonexistence.
The first missing step is an existence argument for a general free
rotund coupled variety. The original conjecture remains
\textbf{UNRESOLVED}.

\subsection{Sources}
\begin{itemize}
\item[S1]F. P. Gallinaro, Exponential sums equations and tropical geometry, Selecta Mathematica (2023).
\url{https://link.springer.com/article/10.1007/s00029-023-00853-y}.
\textit{Formulation source.}
\item[E1]Francesco Gallinaro, \emph{Exponential Sums Equations and
Tropical Geometry}, arXiv:2203.13767v2, author version of [S1].
\url{https://arxiv.org/abs/2203.13767v2}.
\textit{The stated theorem concerns split varieties of the specified
kind; it is not the full coupled-variety assertion. Consulted 27 September 2026.}

\end{itemize}
\end{document}
